\documentclass[10pt]{article}

\usepackage[utf8]{inputenc}

\usepackage[T1]{fontenc}

\usepackage{amsmath}

\usepackage{amsfonts}

\usepackage{amssymb}

\usepackage[version=4]{mhchem}

\usepackage{stmaryrd}

\usepackage{bbold}

\usepackage{graphicx}

\usepackage[export]{adjustbox}

\graphicspath{ {./images/} }



\begin{document}

пию к прежде вычисленным собственным векторам. Приемы исчерпывания могут использоваться и групповыми методами.



Лит.: [1] П а р л е т т Б., Симметричная проблема собственных значений. Численные методы, пер. с англ., М., 1983; [2] У и л к и н с он II ж., Алгебраическая проблема собственных значений, шер. с англ,, М.,

ЧАСТИ. к у р с и в н а я ф у н к ц и, - одно из эквивалентных уточнений понятия вычислимой функции.



ЧАСТИЧНО РЕКУРСИВНЫИ ОПЕРАТОр Е. Плиско. ражение класса всех одноместных функций в себя, определяемос следующим образом. Пусть $\Phi_{z}-$ нек-рый перечисления оператор. С этим оператором естественным образом связан другой оператор $\Psi$, к-рый действует на одноместных функциях. А именно, всякая функция $\varphi$ имеет график - множество всех пар $\langle x, y\rangle$ таких, что $\varphi(x)=y$. При фиксированном способе коди рования пар натуральными числами этот график может рассматриваться как множество $\tau(\varphi)$ натуральных чисел. Если теперь $\Phi_{z}(\tau(\varphi))$ также является графиком нек-рой функции $\psi$, то полагают $\Psi(\varphi)=\psi$. В противном случае считают, что значение $\Psi(\varphi)$ не определено. Таким образом, каждый оператор перечисления $\Phi_{z}$ определяет нек-рый Ч. р. о. $\Psi$.



Если. Ч. р. о. определен на всех функциях, он наз. ре к у р и в ным о і е р а т о р о м. Ч. р. о., к-рый определен на всех всюду определенных функциях и переводит всюду определенные функции во всюду определенные, наз. о б щ е р е к у р с и в ны м о п ер а т о р м. Не всякий Ч. р. о. может быть продолжен до рекурсивного оператора. Всякий общерекурсивный оператор является рекурсивным оператором. Обратное включение не имеет места.



JIum.: [1] Р о д ж е р с X., Теория рекурсивных функций и әффективная вычислимость, пер. с англ., М., 1972.



ЧАСТИЧНО УПоРЯДОЧЕННАЯ ГРУПात па $G$, на к-рой задано отношение частичного порядка $\leqslant$ такое, что для любых $a, b, x, y$ из $G$ неравенство $a \leqslant b$ влетет за собоӥ $x a y \leqslant x b y$.



Множество $P=\{x \in G \mid x \geqslant 1\} \quad$ Ч. у. г., называемое поло жи тельным конусом, или це ло й ч а с т ь ю, группы $G$, обладает свойствами: 1) $P P \subseteq P$; 2) $P \cap P^{-1}=\{1\}$; 3) $x^{-1} P x \subseteq P$ для любых $x \in G$. Всякое подмножество $P$ групшы $G$, удовлетворяющее условиям 1) -3), задает на $G$ частитный порядок $(x \leqslant y$ тогда и только тогда, когда $\left.x^{-1} y \in P\right)$, для к-рого $P$ служит положительным конусом.



Примеры Ч. у. г.: аддитивная группа действительных чисел с обычным порядком; групша $F(X, \mathbb{R})$ функций, заданных на произвольном множестве $X$ со значениями в $\mathbb{R}$, с операцией



\[

(f+g)(x)=f(x)+g(x)

\]



II отношением порядка $f \leqslant g$, если $f(x) \leqslant g(x)$ для всех $x \in X$; группа $A(M)$ всех автоморфизмов линейно упорядоченного множества $M$ на себя, относительно суперпозиции отображений и с отнопением порядка: $\varphi \leqslant \psi$. если $\varphi(m) \leqslant \psi(m)$ для любого $m \in M, \varphi, \psi \in A(M)$.



Основными понятиями теории Ч. у. г. являются понятия порядкового гомоморфизма (см. Упорядоченная әруппа), выпуклой подгруппь, декартова и лексикографич. произведений.



Важнейшие классы Ч. У. г.: линейно упорядоченные әруппы, структурно упорядоченные группы.



Лит.: [1] Б и р к г о ф Г., Теория структур, пер. с англ., M., 1952; [2] Ф у к с Л., Частично упорядоченные алгебраи-



ческие системы, пер. с англ., М., 1965. В. М. Копытов. непустое множество, на к-ром зафиксирован нек-рый порлдок.



Ч. у. м. является гримером модељи. Примеры Ч. у. м.: 1) множество натуральных чисел с обычным порядком; 2) множество натуральных чисел, где $a \leqslant b$ означает, что $a$ делит $b ; 3)$ множество всех подмножеств нек-рого множества, где $a \leqslant b$ означает, что $a \subseteq b$; 4) множество всех действительных функций на отрезке [0, 1], где $f \leqslant g$ означает, что $f(t) \leqslant g(t)$ для всех $t \in[0,1] ; 5)$ множество конечных возрастающих последовательностей натуральных чисел, где



\[

\left(a_{1}, \ldots, a_{k}\right) \leqslant\left(b_{1}, \ldots, b_{l}\right)

\]



означает, что $k \leqslant l$ и $a_{i}=b_{i}$ при $1 \leqslant i \leqslant k$ (см. Дерево); 6) произвольное непустое множество, где $a \leqslant b$ означает $a=b$ (такое Ч. у. м. наз. т р и в и а л ь н ы м, или д и с к р е т ны м).



Каждое Ч. у. м. $P$ можно рассматривать как малую категорию, объектами к-рой служат элементы множества $P$, а множество морфизмов $H(a, b)$ одноэлементно, если $a \leqslant b$, и пусто в остальных случаях. В свою очередь, каждая малая категория, где каждое из множеств $H(a, b)$ содержит не более одного әлемента, определяет 4. у. м.



Если на Ч. у. м. $P$ определить отношение $\triangle$, полагая $a \leq b$, если $b \leqslant a$, то это отношение оказывается порядком. Возникающее таким образом Ч. у. м. наз. д уа льны м, или д вой ствен ны м, к $P$. Отображение $\varphi$ Ч. у. м. $P$ в Ч. у. м. $P^{\prime}$ наз. (анти)изотонным, или (анти)гомоморфизмом, если $a \leqslant b$ в $P$ влечет



\[

\varphi(a) \leqslant \varphi(b) \quad(\varphi(b) \leqslant \varphi(a))

\]



в $P^{\prime}$. Взаимно однозначный (анти)гомоморфизм наз. (анти)изоморфизмом. Тождественное отображение Ч. у. м. $\boldsymbol{P}$ на себя является антиизоморфизмом этого Ч. у. м. на дуальное ему. Изоморфизм является частным случаем резидуального отображения. Последовательное выполнение двух антигомоморфизмов дает гомоморфизм. Совокупность всех Ч. у. м. образует категорию, если морфизмами считать изотонные отображения. Всякое непустое подмножество Ч. У. м. является Ч. У. м. относительно индуцированного на нем порядка.



Если $\boldsymbol{A}$ - непустое подмножество Ч. у. м. $P$, то



\includegraphics[max width=\textwidth, center]{2023_07_09_0a7e30755279506e92a4g-1}

определяется как совокупность всех таких әлементов $x \in P$, что $x \leqslant a(a \leqslant x)$ для всех $a \in A$. Если $a, b \in P$ и $a \leqslant b$, то подмножество



\[

[a, b]=a^{\triangle} \cap b^{\nabla}=\{x \mid a \leqslant x \leqslant b\}

\]



наз. и н т е р в а л о м (хотя с точки зрения словоупотребления, принятого в математич. анализе, здесь следовало бы говорить «отрезок»). Конусы $a^{\triangle}$ и $a \nabla$, где $a \in P$, часто также наз. интервалами. Элемент $u$ из подмножества $A$ наз. н а и б о л ь ІІ и м (н а и м е н ьIII и м), если $a \leqslant u(u \leqslant a)$ для всех $a \in A$. В этом случае $u$ является единственным элементом пересечения $A \cap A \triangle$ $(A \cap A \nabla)$. Наибольший (наименьший) элемент Ч. у. м. $P$ (если он существует) наз. е д и н и ц е й (н у л е м) этого Ч. у. м. Элемент $m$ из подмножества $A$ наз. м а кс и м а л ь ны м (м и н и а ль ны м), если для каждого $x$ из $A$, для к-рого $m \leqslant x(x \leqslant m)$ будет $m=x$. Другими словами,



\[

m^{\triangle} \cap A=m(m \nabla \cap A=m)

\]



Наибольший (наименьший) әлемент является максимальным (минимальным). Обратное верно не всегда. Наименьший (наибольший) әлемент верхнего (нижнего) конуса $A^{\triangle}\left(A^{\nabla}\right)$ наз. т о чно й в е р х в е й (н и ж н е й) г р а н ь ю подмножества $A$ и обозначается как $\sup A$ (inf $A$ ). Расшифровывая это определение, можно сказать, что $u=\sup A$, если $u \geqslant a$ для всех $a \in A$ и $u^{\prime} \geqslant$ $\geqslant u$, если $u^{\prime} \geqslant a$ для всех $a \in A$. Аналогично расшиф-





\end{document}